% Propietario conservado: /Users/ruben/Documents/New project/output/EXTREMA_MEDIA_RAZON_RECIPROCIDAD_ES_EN_20260918/ARTICULO/ES/source/nuclear/07.tex
\chapter{Conservation variationnelle et action orientée}
\label{x_reciprocidad:nuclear:7}
\section{Symétrie de l’action elliptique et charge de Noether}

On fixe une section positive d’action \(\hbar\), la paire angulaire \(A,C^*\) dans une même unité, \(A\ne0\), \(|C^*/A|<1\), et

\[
\eta=\operatorname{artanh}(C^*/A).
\]

Le plan réalisé \(\mathcal P_\eta\) a pour coordonnées \((Q,P)\), toutes deux de dimension racine d’action, avec

\[
\lambda=P\,dQ,\qquad \Omega=dQ\wedge dP=-d\lambda.
\]

La construction de \cref{x_reciprocidad:iv:ellipse:section} donne

\begin{equation}
\mathcal I_\eta(Q,P)=\frac12(e^{-\eta}Q^2+e^\eta P^2),
\qquad H_\eta=\omega\mathcal I_\eta,\qquad \omega>0.
\label{x_reciprocidad:nuc07:eq:1}
\end{equation}

L’ellipse HMT est \(\mathcal I_\eta=\hbar\). Son aire est \(2\pi\hbar\) ; \(\mathcal I_\eta\) a la dimension d’une action et \(H_\eta\) celle d’une énergie. Nous définissons les matrices déjà réalisées

\begin{equation}
M_\eta=\operatorname{diag}(e^{\eta/2},e^{-\eta/2}),\quad
J_2=\begin{pmatrix}0&-1\\1&0\end{pmatrix},\quad
R_\eta(\varepsilon)=M_\eta e^{-\varepsilon J_2}M_\eta^{-1}.
\label{x_reciprocidad:nuc07:eq:2}
\end{equation}

\begin{proposition}
Le groupe \(R_\eta\) est une symétrie canonique de l’action

\begin{equation}
\mathscr S_\eta[z]=\int_{t_a}^{t_b}
\left(P\dot Q-\omega\mathcal I_\eta(Q,P)\right)dt
\label{x_reciprocidad:nuc07:eq:3}
\end{equation}

et sa charge de Noether est \(\mathcal I_\eta\).
\end{proposition}

\begin{proof}
Le déterminant de \(M_\eta\) vaut un et \(M_\eta^{\mathsf T}J_2M_\eta=J_2\). De plus,
\(\mathcal I_\eta(M_\eta y)=\|y\|^2/2\), invariant sous \(e^{-\varepsilon J_2}\).
Le générateur

\[
X_\eta=e^\eta P\,\partial_Q-e^{-\eta}Q\,\partial_P
\]

vérifie \(\iota_{X_\eta}\Omega=d\mathcal I_\eta\). Par la formule de Cartan,

\[
\mathcal L_{X_\eta}\lambda=dF_\eta,\qquad
F_\eta=\lambda(X_\eta)-\mathcal I_\eta=e^\eta P^2-\mathcal I_\eta.
\]

Comme \(X_\eta H_\eta=0\), la variation de la densité d’action est \(dF_\eta/dt\). Sa charge est

\[
\frac{\partial L}{\partial\dot Q}X_\eta Q-F_\eta
=Pe^\eta P-F_\eta=\mathcal I_\eta.
\]

Les équations \(\dot Q=\omega e^\eta P\), \(\dot P=-\omega e^{-\eta}Q\) vérifient directement \(d\mathcal I_\eta/dt=0\).
\end{proof}

Le paramétrage positif de l’ellipse est \(Q=a_S\cos\theta\), \(P=b_S\sin\theta\). Ces équations lui attribuent \(\dot\theta=-\omega\). Par conséquent, au cours d’un tour parcouru selon le flot dynamique, \(\oint P\,dQ=2\pi\mathcal I_\eta\) ; avec \(\theta\) croissant, l’intégrale vaut \(-2\pi\mathcal I_\eta\). C’est la distinction d’orientation établie dans \cref{x_reciprocidad:iv:ellipse:area}. Dans le contour radional, le module vaut \(h=2\pi\hbar\).

\section{Changement de forme : transport et connexion déterminés}

Pour deux formes déjà publiées \(\eta,\eta'\) et une avancée de phase \(\vartheta\), on définit

\begin{equation}
\mathcal T_{\eta',\eta}^{\vartheta}:\mathcal P_\eta\to\mathcal P_{\eta'},
\qquad
\mathcal T_{\eta',\eta}^{\vartheta}
=M_{\eta'}e^{-\vartheta J_2}M_\eta^{-1}.
\label{x_reciprocidad:nuc07:eq:4}
\end{equation}

\begin{proposition}
Ce transport est symplectique et satisfait

\begin{equation}
\mathcal I_{\eta'}(\mathcal T_{\eta',\eta}^{\vartheta}z)
=\mathcal I_\eta(z),\qquad
\mathcal T_{\eta'',\eta'}^{\vartheta_2}
\mathcal T_{\eta',\eta}^{\vartheta_1}
=\mathcal T_{\eta'',\eta}^{\vartheta_1+\vartheta_2}.
\label{x_reciprocidad:nuc07:eq:5}
\end{equation}
\end{proposition}

\begin{proof}
Tous les facteurs sont symplectiques. La première égalité se réduit à l’invariance de \(\|y\|^2/2\) sous une rotation. Dans la seconde, les matrices contiguës \(M_{\eta'}^{-1}M_{\eta'}\) s’annulent et les phases s’additionnent.
\end{proof}

La formule reçoit les formes et les avancées du lecteur ; elle ne sélectionne pas leurs valeurs. La liste complète des transitions peut être retenue avec le transport composé : deux listes donnant la même application composée ne sont pas identifiées comme histoires.

\begin{proposition}
Une réalisation différentiable \(\eta(t),\vartheta(t)\in C^1\), avec \(\dot\vartheta=\omega(t)\), du transport \eqref{x_reciprocidad:nuc07:eq:4} a pour générateur

\begin{equation}
\boxed{
H_{\rm tr}(Q,P,t)=
\omega(t)\mathcal I_{\eta(t)}(Q,P)
+\frac{\dot\eta(t)}2 QP.}
\label{x_reciprocidad:nuc07:eq:6}
\end{equation}

Ce générateur conserve \(\mathcal I_{\eta(t)}(z(t))\).
\end{proposition}

\begin{proof}
En dérivant
\(z(t)=M_{\eta(t)}e^{-\vartheta(t)J_2}y_0\), on obtient

\[
\dot Q=\omega e^\eta P+\frac{\dot\eta}2Q,\qquad
\dot P=-\omega e^{-\eta}Q-\frac{\dot\eta}2P.
\]

Ce sont les équations de Hamilton de \eqref{x_reciprocidad:nuc07:eq:6}. Le terme de connexion réalise exactement \(\dot M_\eta M_\eta^{-1}z\), et son Hamiltonien est déterminé à une fonction de \(t\) seul près. En effet,

\begin{equation}
\begin{aligned}
\frac{d\mathcal I_\eta}{dt}
&=\frac{\dot\eta}{2}(-e^{-\eta}Q^2+e^\eta P^2)
+\left\{\mathcal I_\eta,\frac{\dot\eta}{2}QP\right\}\\
&=\frac{\dot\eta}{2}(-e^{-\eta}Q^2+e^\eta P^2)
+\frac{\dot\eta}{2}(e^{-\eta}Q^2-e^\eta P^2)=0.
\end{aligned}
\label{x_reciprocidad:nuc07:eq:7}
\end{equation}

De manière équivalente, \(Q=e^{\eta/2}q\), \(P=e^{-\eta/2}p\) donnent

\begin{equation}
P\dot Q-H_{\rm tr}
=p\dot q-\frac{\omega(t)}2(q^2+p^2).
\label{x_reciprocidad:nuc07:eq:8}
\end{equation}

La symétrie de rotation et sa charge sont transportées exactement.
\end{proof}

Ici, le terme de connexion a été dérivé du transport, non d’une force extérieure. L’hamiltonien LC de forme fixe conserve \(\mathcal I_\eta\). Lors d’une déformation, omettre le second terme de \eqref{x_reciprocidad:nuc07:eq:6} laisse le premier terme de \eqref{x_reciprocidad:nuc07:eq:7} non annulé. La conservation variationnelle exige de transporter également la forme canonique.

La proposition détermine le générateur de l’interpolation déclarée ; elle n’attribue pas une interpolation physique unique à toute histoire TPK. L’identité discrète \eqref{x_reciprocidad:nuc07:eq:4}–\eqref{x_reciprocidad:nuc07:eq:5} fonctionne sans interpolation.

\section{Orientation et bord de l’action}

Les transformations de \cref{x_reciprocidad:iv:ellipse:reciprocidad} distinguent \(S(Q,P)=(P,Q)\) et \(J_2(Q,P)=(-P,Q)\). Toutes deux échangent la forme \(\eta\leftrightarrow-\eta\), mais

\begin{equation}
S^*\Omega=-\Omega,\quad J_2^*\Omega=\Omega,\qquad
S^*\lambda=-\lambda+d(PQ),\quad
J_2^*\lambda=\lambda-d(PQ).
\label{x_reciprocidad:nuc07:eq:9}
\end{equation}

De même,

\begin{equation}
S R_\eta(\varepsilon)S^{-1}=R_{-\eta}(-\varepsilon),\qquad
J_2R_\eta(\varepsilon)J_2^{-1}=R_{-\eta}(\varepsilon).
\label{x_reciprocidad:nuc07:eq:10}
\end{equation}

Elles se vérifient en utilisant \(SM_\eta=M_{-\eta}S\), \(J_2M_\eta=M_{-\eta}J_2\), \(SJ_2S^{-1}=-J_2\). Pour une courbe ouverte orientée \(z:[a,b]\to\mathcal P_\eta\),

\begin{equation}
\int_{J_2z}\lambda=\int_z\lambda-[PQ]_a^b,\qquad
\int_{Sz}\lambda=-\int_z\lambda+[PQ]_a^b.
\label{x_reciprocidad:nuc07:eq:11}
\end{equation}

Les termes de bord s’annulent dans un cycle et se compensent lors de la concaténation d’extrémités compatibles. L’énergie quadratique seule ne détermine ni ces signes ni ces termes.

La monodromie nonadique enregistre le retour de phase et \(q_{\rm mem}\mapsto q_{\rm mem}+1\). La constance d’une charge n’impose pas au compteur de rester fixe. Dans la réalisation elliptique, un tour conserve \(\mathcal I_\eta\) tout en augmentant l’intégrale cumulée orientée de \(2\pi\mathcal I_\eta\). La charge, l’action intégrée, la phase et le registre ont des types distincts. On n’identifie pas ici un tour angulaire à un cycle TPK quelconque sans spécifier son lecteur.

\section{Charge variationnelle de l’horloge et bilan entre lecture et mémoire}

\subsection{Une charge d’action explicite de l’horloge}

Le calendrier se réalise sur \(\mathcal H_{\rm clk}=\mathbb C^{108}\), avec les projecteurs de Fourier \(P_k=|\widetilde k\rangle\langle\widetilde k|\), et détermine

\begin{equation}
N_{108}=\sum_{k=0}^{107}kP_k,\quad
H_{\rm clk}=\hbar\omega_0N_{108},\quad
\omega_0=\frac{2\pi}{108t_0}.
\label{x_reciprocidad:nuc07:eq:12}
\end{equation}

Son propagateur pendant \(t_0\) est le déplacement d’une cellule. L’action de l’horloge est

\begin{equation}
\mathscr S[\psi]=\int
\left[\frac{i\hbar}{2}
(\langle\psi,\dot\psi\rangle-\langle\dot\psi,\psi\rangle)
-\langle\psi,H_{\rm clk}\psi\rangle\right]dt.
\label{x_reciprocidad:nuc07:eq:13}
\end{equation}

Le produit est antilinéaire en la première variable. L’équation est \(i\hbar\dot\psi=H_{\rm clk}\psi\).

\begin{proposition}
Pour \(B=B^*\) avec \([B,H_{\rm clk}]=0\), la symétrie \(\psi\mapsto e^{-i\varepsilon B}\psi\) a pour charge

\begin{equation}
\mathcal Q_B(\psi)=\hbar\langle\psi,B\psi\rangle.
\label{x_reciprocidad:nuc07:eq:14}
\end{equation}
\end{proposition}

\begin{proof}
En localisant le paramètre, \(\delta\psi=-i\varepsilon(t)B\psi\), on obtient
\(\delta L=\hbar\dot\varepsilon\langle\psi,B\psi\rangle\).
L’intégration par parties dans les équations variationnelles donne \(d\mathcal Q_B/dt=0\). Directement,

\[
\frac{d}{dt}\langle\psi,B\psi\rangle
=\frac{i}{\hbar}\langle\psi,[H_{\rm clk},B]\psi\rangle=0.
\]

Avec \(B=I\), la norme est conservée ; avec \(B=N_{108}\), la charge est l’action \(E_{\rm clk}/\omega_0\).
\end{proof}

La forme symplectique de l’espace réel sous-jacent est
\(\Omega_{\rm clk}(u,v)=2\hbar\operatorname{Im}\langle u,v\rangle\) ;
sa forme de degré un est
\(\Theta_\psi(v)=-\hbar\operatorname{Im}\langle\psi,v\rangle\),
avec \(\Omega_{\rm clk}=-d\Theta\).

\subsection{Bilan exact de charge}

Soient \(\mathcal H,\mathcal K\) des espaces hermitiens finis, \(J:\mathcal H\to\mathcal K\) une isométrie, \(U\) unitaire sur \(\mathcal K\), et

\begin{equation}
\Pi=JJ^*,\quad T=J^*UJ,\quad
\eta_{\rm mem}=(I-\Pi)UJ.
\label{x_reciprocidad:nuc07:eq:15}
\end{equation}

L’opérateur \(\eta_{\rm mem}\) est distinct du paramètre scalaire \(\eta\).

\begin{proposition}
Si

\begin{equation}
\widetilde B=\widetilde B^*,\quad U^*\widetilde BU=\widetilde B,
\quad [\widetilde B,\Pi]=0,\quad B=J^*\widetilde BJ,
\label{x_reciprocidad:nuc07:eq:16}
\end{equation}

alors

\begin{equation}
B=T^*BT+\eta_{\rm mem}^*\widetilde B\eta_{\rm mem},
\qquad
\mathcal Q_B(\psi)=
\mathcal Q_B(T\psi)+\mathcal Q_{\widetilde B}(\eta_{\rm mem}\psi).
\label{x_reciprocidad:nuc07:eq:17}
\end{equation}
\end{proposition}

\begin{proof}
On développe \(UJ\psi=JT\psi+\eta_{\rm mem}\psi\) en
\(J^*U^*\widetilde BUJ=J^*\widetilde BJ\).
La commutation avec \(\Pi\) annule les termes croisés. Évaluer l’identité opératorielle et multiplier par \(\hbar\) donne le bilan d’action.
\end{proof}

Par itération,

\begin{equation}
\mathcal Q_B(\psi)=\mathcal Q_B(T^n\psi)
+\sum_{j=0}^{n-1}\mathcal Q_{\widetilde B}
(\eta_{\rm mem}T^j\psi).
\label{x_reciprocidad:nuc07:eq:18}
\end{equation}

Cette somme enregistre les compléments successifs de la dynamique comprimée. La dynamique complète a une autre formule exacte, avec \(z_n=U^nJ\psi\) :

\begin{equation}
\mathcal Q_B(\psi)=
\mathcal Q_B(J^*z_n)+
\mathcal Q_{\widetilde B}((I-\Pi)z_n).
\label{x_reciprocidad:nuc07:eq:19}
\end{equation}

On n’identifie pas \(T^n\) à \(J^*U^nJ\) sans autre propriété : l’état complet peut réabsorber de la mémoire. Si \(B,\widetilde B\ge0\), tous les termes de \eqref{x_reciprocidad:nuc07:eq:18} sont non négatifs ; pour des charges signées, l’égalité demeure, sans cette interprétation de positivité.

Si \(U\) est le propagateur d’une action comme \eqref{x_reciprocidad:nuc07:eq:13}, la commutation de \(\widetilde B\) avec son Hamiltonien apporte la conservation de Noether. La condition de réduction par \(\Pi\) spécifie les lectures qui séparent la charge sans termes d’interférence.

Sans imposer \eqref{x_reciprocidad:nuc07:eq:16}, l’identité générale conserve les termes supplémentaires :

\begin{equation}
\begin{aligned}
&\mathcal Q_B(T\psi)+\mathcal Q_{\widetilde B}(\eta_{\rm mem}\psi)
+2\hbar\operatorname{Re}
\langle JT\psi,\widetilde B\eta_{\rm mem}\psi\rangle\\
&=\mathcal Q_B(\psi)+
\hbar\langle\psi,J^*(U^*\widetilde BU-\widetilde B)J\psi\rangle.
\end{aligned}
\label{x_reciprocidad:nuc07:eq:20}
\end{equation}

La charge physique utilise son lecteur, par exemple \(N_{108}\) ; une identité de normes correspond à \(B=I\). Les sections \(\hbar_{\rm pre}\) et \(\hbar_{\rm ret}\) conservent le lecteur de retour de \cref{x_reciprocidad:iv:action:secciones,x_reciprocidad:iv:action:razon} ; la différence entre elles et le complément de mémoire possèdent leurs domaines de définition respectifs.

\section{Bilan covariant du tenseur total}

L’élimination variationnelle de la contorsion et l’identité de Noether de l’action covariante, développées dans la proposition~\ref{x_reciprocidad:vii:prop:fuente-metrica-balance}, fournissent

\begin{equation}
\mathcal T_{\mu\nu}
=T^{\rm phys,LC}_{\mu\nu}+U^{\rm phys,spin}_{\mu\nu},
\qquad
\mathring\nabla^\mu\mathcal T_{\mu\nu}=0.
\label{x_reciprocidad:nuc07:eq:21}
\end{equation}

Son domaine est la carte HMT–MD réalisée, avec les couplages maintenus constants pendant la variation. Pour un champ de Killing \(\xi\),

\begin{equation}
j^\mu_\xi=\mathcal T^{\mu\nu}\xi_\nu,\qquad
\mathring\nabla_\mu j^\mu_\xi=0.
\label{x_reciprocidad:nuc07:eq:22}
\end{equation}

La preuve utilise \eqref{x_reciprocidad:nuc07:eq:21}, la symétrie de \(\mathcal T\) et l’antisymétrie de \(\mathring\nabla_\mu\xi_\nu\). Après intégration, la différence entre les charges des hypersurfaces est le flux à travers la frontière latérale.

La conservation appartient au tenseur total ; chaque composante peut échanger de la charge. Le théorème~\ref{x_reciprocidad:vii:thm:intercambio-traza-residual} détermine explicitement \(\mathcal J^{\rm tr}_\nu=\partial_\nu\operatorname{tr}_g(T^{\rm res})/4\) et les bilans opposés de la trace et de la partie de trace nulle. Cela est retrouvé comme une instance covariante du critère de bilan conjoint, non comme une nouvelle équation gravitationnelle.
